
\subsubsection{Reward Model Evaluation}

RewardBench \citep{lambert2024rewardbench} established the first standardized benchmark for evaluating reward models, measuring accuracy across chat, safety, and reasoning categories. Top models achieve 80--85\% aggregate accuracy. This aggregate view treats all errors equally and provides no decomposition of where failures occur. The present work complements RewardBench by introducing mode-based decomposition that stratifies performance by human-RM agreement patterns.

The Alignment Ceiling \citep{lambert2023alignmentceiling} identified objective mismatch in RLHF: reward models overoptimize proxy objectives that diverge from true human intent. This work demonstrated that reward model training contains systematic biases but did not quantify the prevalence of overconfident misalignment in preference data. The present work operationalizes ''overconfident misalignment'' as Mode 3 and measures its proportion.

\subsubsection{LLM-as-Judge Evaluation}

Wei et al. (2024) systematically evaluated LLM judges for alignment tasks, finding that prompt templates and judge models significantly affect reliability. The present work differs by examining reward model confidence rather than LLM judge outputs, and by conditioning analysis on human disagreement rather than treating human labels as ground truth.

\subsubsection{Bidirectional Alignment Frameworks}

Shen et al. (2024) proposed the Bidirectional Human-AI Alignment framework, systematically reviewing over 400 papers spanning ML, HCI, and NLP. The present work contributes to AI-to-human alignment measurement by revealing that reward models exhibit systematic overconfidence on human-disagreement samples.
